\chapter{Testes e validação}

Este capítulo apresenta a estratégia de testes, as suítes que compõem a validação automatizada, os
testes de arquitetura que verificam os princípios de projeto, o processo de integração contínua e os
procedimentos de validação manual.

\section{Estratégia}

A estratégia segue a pirâmide de testes \cite{cohn2009}. A base é formada por testes unitários rápidos
das regras de domínio, das políticas de aplicação e do processamento de sinais. O nível intermediário
reúne testes de integração dos adaptadores com dublês das bibliotecas externas e do banco SQLite real
em arquivo temporário. No topo ficam poucos testes que abrem a janela Tk real, marcados como lentos e
executados com servidor gráfico virtual na integração contínua.

Três decisões de projeto tornaram essa estratégia viável. A primeira é a arquitetura hexagonal: todo
serviço externo (modelo de linguagem, reconhecedores, sintetizador, reprodutor, lançador de programas)
é acessado por uma porta e pode ser substituído por um dublê \cite{meszaros2007}. A segunda é a injeção
do relógio, que torna determinísticos os testes de lembretes e de períodos. A terceira é a
reprodução simulada com duração configurável, que permite provar propriedades temporais do pipeline,
como a de que a primeira frase começa a tocar antes de o modelo terminar a resposta.

\section{Dublês de teste}

O módulo comum de testes oferece dublês para cada porta:
\begin{enumerate}[label=\alph*)]
  \item modelos roteirizados, que devolvem pedaços de texto com atraso controlado;
  \item um cliente falso do modelo Claude, que reproduz o protocolo de transmissão em fluxo contínuo com
        blocos de texto e de ferramenta;
  \item sintetizadores que registram os textos e podem falhar sob demanda;
  \item reprodutores que registram as faixas tocadas e acionam ganchos durante a reprodução;
  \item reconhecedores com falas roteirizadas e transcrições opcionais;
  \item detectores de ativação que devolvem ativações ou falhas roteirizadas e depois aguardam o
        cancelamento;
  \item lançadores de programas vazios ou com listas controladas;
  \item um servidor HTTP falso do Ollama e catálogos de modelos roteirizados, que respondem com listas
        ou com falhas (sem chave, sem conexão, servidor fora do ar);
  \item listas de microfones e estados de inicialização controlados, sem depender do hardware.
\end{enumerate}
Uma fixture automática aponta a pasta pessoal e a pasta de configuração para diretórios temporários
em todos os testes, de modo que nenhum deles possa criar atalhos de inicialização ou comandos na pasta
real do usuário.
Também são oferecidos geradores de áudio sintético, como senoides em frequências conhecidas e arquivos
WAV temporários, para validar a análise de tom e de envelope.

\section{Suítes de teste}

O Quadro~\ref{qua:suites} descreve as suítes e o número de casos de teste de cada uma, contando as
variações parametrizadas.

\begin{singlespace}
\legendaquadro{Suítes de teste automatizado}\label{qua:suites}
\small
\begin{tabularx}{\textwidth}{|P{4.2cm}|C{1.3cm}|L|}
\hline
\textbf{Suíte} & \textbf{Casos} & \textbf{O que garante} \\ \hline
\endfirsthead
\hline
\textbf{Suíte} & \textbf{Casos} & \textbf{O que garante} \\ \hline
\endhead
\cod{test_i18n} & 190 & Catálogo completo e com os mesmos campos nas duas línguas; comandos, períodos e respostas em inglês; troca ao vivo; precedência da configuração \\ \hline
\cod{test_ativacao} & 79 & Gramática de “ok Fael” e busca fonética com saídas reais do reconhecedor local; frases parecidas que não podem disparar; extração do comando com saídas reais do reconhecedor remoto; vigia liberando o microfone; desligamento após falhas \\ \hline
\cod{test_memoria} & 41 & Dezenove expressões de período, intenção “o que eu fiz”, resultado líquido e respostas \\ \hline
\cod{test_programas} & 29 & Leitura de entradas \cod{.desktop}, extração do pedido, ranqueamento, sinônimos, ambiguidade, abertura sem interpretador de comandos e prioridade da agenda \\ \hline
\cod{test_ollama} & 24 & Respostas em camadas, ferramentas seguras, chamadas escritas como texto, queda para o modo offline, interrupção, histórico, idioma, cliente HTTP contra servidor falso e escolha do cérebro \\ \hline
\cod{test_sqlite} & 18 & Esquema versionado, transação conjunta de estado e histórico, gatilhos de somente acréscimo, detecção de quatro tipos de adulteração, consulta por período, concorrência, migração e preferências \\ \hline
\cod{test_agenda} & 18 & Regras da agenda e transações (persiste e notifica somente em caso de sucesso) \\ \hline
\cod{test_analise_audio} & 23 & Estimativa de tom com erro de até 2\%, envelope de volume e timbre e corte do silêncio das pontas (sem cortar consoantes fracas) \\ \hline
\cod{test_orquestrador} & 17 & Prova do pipeline, interrupção, modo contínuo, anúncios em fila, falhas de modelo, síntese e microfone, corrida do cancelamento e pedidos digitados durante a fala ou a escuta \\ \hline
\cod{test_modelo_offline} & 15 & Comandos, plural, lembretes relativos e despedida \\ \hline
\cod{test_animacao} & 15 & Suavização independente da taxa de quadros e seleção do quadro da boca \\ \hline
\cod{test_modelo_claude} & 14 & Transmissão real do texto, laço de ferramentas, limite de tamanho, recusa, erros da API, descarte da rodada e limite do histórico \\ \hline
\cod{test_arquitetura} & 15 & Regra de dependência entre camadas e princípios SOLID (Seção~\ref{sec:testes-arquitetura}) \\ \hline
\cod{test_configuracoes} & 15 & Conversões de tom e ritmo, precedência das preferências, janela de tamanho fixo e diálogo de configurações, inclusive microfone, inicialização, modelo em uso, mudança e backup do banco \\ \hline
\cod{test_microfone_inicializacao} & 14 & Lista de microfones sem monitores, variável de escolha restrita à abertura, rótulos, reabertura da escuta passiva, atalho de sessão e comandos no terminal sem apagar arquivos de outra origem \\ \hline
\cod{test_banco_e_modelos} & 13 & Backup conferido, mudança do banco com gravações antes e depois, recusa de sobrescrever, banco ausente na abertura, catálogos do Claude e do Ollama e texto da lista de modelos \\ \hline
\cod{test_perfil_voz} & 12 & Criação, leitura e tolerância a arquivos inexistentes ou corrompidos \\ \hline
\cod{test_ferramentas} & 9 & Esquemas e validação de entradas (tipos, campos extras, booleano diferente de inteiro) \\ \hline
\cod{test_avatar} & 9 & Busto centralizado, camisa, tamanho uniforme dos quadros e cache \\ \hline
\cod{test_vad}, \cod{test_estados_barramento}, \cod{test_frases} & 23 & Detecção de fala, transições de estado, entrega de eventos entre threads e divisão de frases \\ \hline
\cod{test_sessao_pipeline_vigia} & 8 & Pipeline e vigia isolados, síntese paralela com ordem preservada e aplicação das preferências \\ \hline
\cod{test_renderizacao} & 8 & Canvas diferencial, fila de trabalho com orçamento, laço de quadros adaptativo, halos opacos e regressão de desempenho do palco em repouso \\ \hline
\cod{test_ganho}, \cod{test_configuracao_tema}, \cod{test_adaptadores_voz} & 21 & Ganho automático, configuração (inclusive os nomes antigos das variáveis) e tema, cache LRU de voz e reprodutor com mixer real \\ \hline
\cod{test_agendador}, \cod{test_repositorio_json}, \cod{test_bootstrap_ui} & 13 & Lembrete no horário exato, formato legado, composição, queda para o modo offline e janela Tk real \\ \hline
\textbf{Total} & \textbf{643} & \\ \hline
\end{tabularx}
\fonte{elaboração própria (2026), a partir da coleta do pytest.}
\end{singlespace}

\section{Testes de arquitetura}\label{sec:testes-arquitetura}

Inspirados nas funções de aptidão arquitetural de \textcite{ford2017}, os testes de arquitetura
transformam as regras de projeto em verificações executáveis. Eles estão no arquivo
\cod{tests/test_arquitetura.py} e verificam:

\begin{enumerate}[label=\alph*)]
  \item \textbf{regra de dependência}: a árvore sintática de cada módulo é analisada, todas as
        importações internas são extraídas e cada uma é comparada com uma tabela que declara, para cada
        camada, as camadas que ela pode importar; qualquer violação é relatada com arquivo, linha e
        sentido da dependência;
  \item \textbf{cobertura das regras}: toda camada existente no pacote precisa ter uma regra declarada,
        para que uma camada nova não escape da verificação;
  \item \textbf{pureza do núcleo}: o domínio e as portas não podem importar bibliotecas de
        infraestrutura (banco de dados, interface gráfica, processamento numérico, áudio, rede, modelo de
        linguagem nem execução de processos);
  \item \textbf{LSP}: o banco SQLite precisa oferecer todos os métodos de todos os contratos de
        persistência, e o repositório legado não pode ter a assinatura que promete guardar o histórico;
  \item \textbf{DIP}: a agenda precisa funcionar com um repositório em memória escrito no próprio teste,
        sem banco de dados;
  \item \textbf{OCP}: toda ferramenta declarada ao modelo precisa ter uma ação registrada e vice-versa, e
        uma intenção nova acrescentada à cadeia do modo offline precisa funcionar sem alterar o código
        existente;
  \item \textbf{SRP}: os métodos do pipeline e da vigia não podem voltar ao orquestrador, e as políticas
        da sessão não podem voltar à raiz de composição.
\end{enumerate}

\section{Análise estática e integração contínua}

O comando \cod{make verificar} executa, em sequência, a análise estática (ruff), a verificação de tipos
(mypy) e a suíte completa (pytest). O mesmo conjunto é executado pelos ganchos de \en{pre-commit} e
pela integração contínua no GitHub Actions, em Python 3.10 e 3.12, com o servidor gráfico virtual Xvfb
para os testes de interface. A medição de cobertura falha se o total ficar abaixo de 80\%.

\section{Cobertura}

A cobertura combinada de linhas e ramificações foi de 86,6\% sobre 4.694 comandos. A
Tabela~\ref{tab:cobertura} apresenta a cobertura por camada.

\begin{singlespace}
\legendatabela{Cobertura de testes por camada (linhas e ramificações)}\label{tab:cobertura}
\small
\begin{tabularx}{\textwidth}{LC{2.6cm}C{3.2cm}}
\toprule
\textbf{Camada} & \textbf{Comandos} & \textbf{Cobertura (\%)} \\ \midrule
\endfirsthead
\toprule
\textbf{Camada} & \textbf{Comandos} & \textbf{Cobertura (\%)} \\ \midrule
\endhead
Portas & 85 & 100,0 \\
Internacionalização & 68 & 98,7 \\
Domínio & 551 & 98,2 \\
Aplicação & 758 & 95,5 \\
Avatar & 290 & 95,2 \\
Configuração & 117 & 93,8 \\
Áudio & 260 & 90,8 \\
Adaptadores & 1.297 & 82,1 \\
Apresentação & 1.130 & 76,0 \\
Raiz de composição & 136 & 75,0 \\ \midrule
\textbf{Total} & \textbf{4.694} & \textbf{86,6} \\
\bottomrule
\end{tabularx}
\fonte{elaboração própria (2026), com dados do pytest-cov.}
\end{singlespace}

A distribuição é coerente com a arquitetura: as camadas internas, que concentram as regras, têm
cobertura próxima de 100\%. As menores coberturas estão nos adaptadores de hardware (microfone,
áudio nativo e escuta local), cuja execução exige dispositivos reais, e na interface gráfica, cujos
caminhos de desenho são parcialmente exercitados pelos testes com a janela real.

\section{Validação manual}

Os aspectos que dependem de hardware ou de percepção humana foram validados manualmente:
\begin{enumerate}[label=\alph*)]
  \item a aparência do avatar, por meio de prévias geradas com as mesmas imagens e posições da interface;
  \item o reconhecimento da frase de ativação com a voz real do usuário, acompanhado pelos registros
        detalhados, que mostram o texto reconhecido, o volume de pico e o ganho aplicado;
  \item a gravação real da voz e a criação do perfil;
  \item a migração dos dados reais, feita primeiro em uma cópia;
  \item a execução em inglês;
  \item a abertura da janela a partir da instalação simulada do pacote, com a verificação da classe da
        janela e da criação das pastas de dados e de cache do usuário;
  \item a escolha do microfone, conferida na lista de gravações ativas do servidor de áudio, que mostrou
        o assistente gravando do microfone escolhido, e não do padrão;
  \item a tela de configurações com os serviços reais e uma pasta de dados temporária, mostrando o
        modelo local em uso, os modelos instalados e um backup conferido;
  \item o atalho de sessão gerado, validado pela ferramenta \cod{desktop-file-validate}.
\end{enumerate}

\section{Estabilidade da suíte}

Como o sistema é concorrente, uma suíte aprovada uma única vez não basta. Antes de cada entrega, a suíte
completa foi executada repetidamente, inclusive com a máquina sob carga. Esse procedimento revelou e
permitiu corrigir quatro defeitos intermitentes ao longo do projeto: dois testes que dependiam de
tempo sob carga, a liberação de variáveis da interface fora da thread principal e o bloqueio da thread
de renderização pelo coletor de lixo, descrito no Capítulo 6.
